\documentclass[../../main.tex]{subfiles}
\begin{document}

\begin{center}
  {\fontsize{48pt}{58pt}\selectfont \textbf{1976}\normalsize\textbf{年度}}
\end{center}

\begin{table}[htb]
  \centering
  \caption{出題テーマと難易度}
  \begin{tabular}{|c|p{9cm}|c|}
    \hline
    \textbf{問題番号} & \textbf{テーマ} & \textbf{難易度} \\
    \hline
    第1問 & 多項式の区間上の大きさと値の制御 & \\
    \hline
    第2問 & 内積が整数条件を満たすベクトルの角と長さ & \\
    \hline
    第3問 & 微分と積分を含む多項式の恒等式 & \\
    \hline
    第4問 & 不等式で制限された数列の極限 & \\
    \hline
    第5問 & 微分可能な関数の局所的な切り替え & \\
    \hline
    第6問 & 袋の判定を誤る確率の極限 & \\
    \hline
  \end{tabular}
\end{table}

\begin{table}[htb]
  \centering
  \caption{解答}
  \renewcommand{\arraystretch}{1.4}
  \begin{tabular}{|c|c|p{9cm}|}
    \hline
    \textbf{問題番号} & \textbf{小問} & \textbf{答え} \\
    \hline
    第1問 & (1) & $n\ge10$ \\
    \cline{2-3}
    & (2) & 例えば$g(x)=(x-\dfrac{1}{2})^{12}$ \\
    \hline
    第2問 & (i) & 角は$0^\circ,30^\circ,45^\circ,60^\circ,90^\circ$またはこれらの補角 \\
    \cline{2-3}
    & (ii) & $0^\circ$では長さ比$1:1$または$1:2$，$30^\circ$では$1:\sqrt{3}$，$60^\circ$では$1:1$ \\
    \cline{2-3}
    & (iii) & 条件を満たす$12$本のベクトルの例は解答を参照 \\
    \hline
    第3問 & -- & $f(x)=\dfrac{4}{9}$，$-\dfrac{x^2}{6}+x-\dfrac{5}{6}$，$-\dfrac{x^2}{6}\pm\dfrac{x}{\sqrt{2}}+\dfrac{1}{6}$ \\
    \hline
    第4問 & (i) & $\displaystyle\lim_{n\to\infty}a_n=1$ \\
    \cline{2-3}
    & (ii) & $(a_n-1)^2<\dfrac{1}{4n}$ \\
    \hline
    第5問 & (i) & $h(x)=g(x)f(|x|/a)$ \\
    \cline{2-3}
    & (ii) & 例えば$[0,\dfrac{1}{2}]$で$2x^2$，$[\dfrac{1}{2},1]$で$1-2(x-1)^2$とし，外側では$0,1$とする \\
    \hline
    第6問 & -- & $\dfrac{14}{75}$ \\
    \hline
  \end{tabular}
\end{table}

\paragraph{試験形態}
6問．試験時間は要確認．

\end{document}
